\section{Summary}
Computational-physics PhD researcher who turns many-body and first-principles models into scalable scientific software. Combines Python/C++, MPI/OpenMP, CUDA/ROCm, PETSc/SLEPc, and scientific ML for materials calculations on leadership-class GPU systems.

\section{Technical Skills}
\SkillLine{Scientific computing}{Python, C++, MPI, OpenMP, PETSc, SLEPc, mpi4py, petsc4py, slepc4py, DOLFINx, NumPy, SciPy, SymPy}
\SkillLine{GPU and HPC}{CUDA (cuBLAS, cuTENSOR), ROCm (rocBLAS, rocFFT), CUDA-aware MPI, collective communication, strong/weak scaling, CMake, TotalView, Linaro DDT; Frontier (AMD MI250X), Perlmutter (NVIDIA A100), Frontera}
\SkillLine{ML and systems}{PyTorch, PyTorch Geometric, JAX, equivariant GNNs, automatic differentiation, Docker, Kubernetes, Linux}

\section{Experience}
\Role{PhD Researcher}{Yale University}{Electrical Engineering}{2021--present}
\begin{itemize}
\item Formulate and implement first-principles many-body models of self-trapped excitons and exciton-polarons in LiF and lead-free double perovskites.
\item Maintain distributed Python/C++ simulation workflows using MPI, OpenMP, PETSc, and SLEPc for INCITE-supported workloads on Frontier and Perlmutter.
\item Port multi-GPU numerical workloads across CUDA and ROCm; use scaling studies, collective communication, and TotalView/Linaro DDT to isolate distributed failures.
\item Construct DeePH-style computational graphs and train equivariant graph neural networks with PyTorch Geometric and automatic differentiation to accelerate optical-property calculations; parts of the workflow achieved roughly 100x speedup by avoiding separate expensive derivative calculations.
\item Co-author a 2025 \textit{Journal of the American Chemical Society} paper and present related first-principles research at APS meetings.
\end{itemize}
\Role{Software and Embedded Systems Engineer}{Probe Technologies / Weatherford Wireline Services}{Houston, TX}{2019--2021}
\begin{itemize}
\item Implemented CUDA/OpenCL waveform-processing and OpenGL visualization tools for acoustic well-logging diagnostics.
\item Delivered C\# and ASP.NET services backed by Microsoft SQL Server for remote tool testing, sensor calibration, and engineering data management.
\end{itemize}

\section{Selected Projects}
\Project{Finite-element simulation portfolio}{Implemented transient heat, Poisson, and incompressible-flow cases in DOLFINx/PETSc, with PyVista and Matplotlib post-processing.}
\Project{EGNN force model}{Trained an equivariant molecular model on QM9 using PyTorch Geometric/e3nn; differentiated predicted energies to obtain forces.}
\Project{H2 VQE benchmark}{Built a reproducible PySCF-to-Qiskit workflow with Jordan--Wigner mapping, UCCSD VQE, and exact-diagonalization reference curves.}

\section{Education}
\Role{PhD, Electrical Engineering}{Yale University}{Computational materials, scientific computing, and ML}{2021--present}
\Role{MPhil and MS, Electrical Engineering}{Yale University}{}{2023}
\Role{MEng and BS, Electrical and Computer Engineering}{Cornell University}{}{2013--2018}

\ifcv
\section{Selected Publications and Presentations}
\textbf{Journal article:} de Paula et al., including S. Vadivel, ``Time-domain observation of ultrafast self-trapped exciton formation in lead-free double halide perovskites,'' \textit{JACS}, 2025.\\
\textbf{APS presentations:} First-principles studies of self-trapped excitons and exciton-polarons, 2024--2026.\\
\section{Additional Technical Work}
\Project{Containerized scientific software}{Packaged CUDA-ready and cross-environment first-principles stacks with Docker and published reusable images to Docker Hub.}
\fi
